\documentclass[a4paper,12pt]{article}

\input{tools/packages.tex}
\input{tools/styles.tex}
\input{tools/acro.tex}

\setcounter{secnumdepth}{4}

\titleformat{\paragraph}
{\normalfont\normalsize\bfseries}{\theparagraph}{1em}{}
\titlespacing*{\paragraph}
{0pt}{3.25ex plus 1ex minus .2ex}{1.5ex plus .2ex}
\begin{document}

%%% COVER %%%
\fancypagestyle{alim}{\fancyhf{}\renewcommand{\headrulewidth}{0pt}
  \cfoot{
\includegraphics[height=2.2cm]{img/logos/logo_telecos.png}}
}
\thispagestyle{empty}
\begin{center}
  {\sffamily
  \resizebox{0.8\textwidth}{!}{
\includegraphics{img/logos/upc_completo+telecos.png}}\\
  \vspace{1cm}
  {\Huge Work-flux optimization for the use of GRASP algorithm in quasi-real-time}\\
  \vspace{0.5cm}
  {\color{black}\hrule height 1pt}
  \vspace{1cm}
  {\large{Master Thesis\\
      submitted to the Faculty of the \\
      Escola T\`ecnica d'Enginyeria de Telecomunicaci\'o de Barcelona \\
      Universitat Polit\`ecnica de Catalunya \\
      by\\
      \vspace{0.4cm}
      Mireia López Barrón}}

  \vspace{1.5cm}

  {In partial fulfillment\\
    of the requirements for the master in\\
    \textit{(TELECOMMUNICATIONS)} \textbf{ENGINEERING}}

  \vspace{2cm}

  {{Advisor: Alejandro Rodriguez Gomez}} \\
  {{Barcelona, Date XXXXX}}
  \thispagestyle{alim}
  }

  %%% INDEX %%%
\end{center}
\newpage
\tableofcontents

%%% LISTS %%%
\newpage
\listoffigures
\listoftables

%%% REVISION %%%
\newpage
\section*{Revision history and approval record}

\selectlanguage{english}
%\foreignlanguage{english}

\begin{center}

\begin{tabularx}{\textwidth}{|p{2cm}|p{2.5cm}|X|}
\hline
\textbf{Revision} &
\textbf{Date} &
\textbf{Purpose}\\\hline
0 &
10/12/2025 &
Document creation\\\hline
1 &
13/01/2026 &
Document revision\\\hline
~
 &
~
 &
~
\\\hline
~
 &
~
 &
~
\\\hline
~
 &
~
 &
~
\\\hline
\end{tabularx}
\end{center}



\bigskip

{\selectlanguage{english}
DOCUMENT DISTRIBUTION LIST}

\begin{center}
\tablefirsthead{}
\tablehead{}
\tabletail{}
\tablelasttail{}
\begin{tabularx}{\textwidth}{|X|X|}
\hline
\textbf{Name} &
\textbf{e-mail}\\\hline
{Mireia L\'opez Barr\'on} &
~ mireia.lopez.barron\@estudiantat.upc.edu
\\\hline
{Alejandro Rodriguez Gomez} &
~ alejandro.rodriguez.gomez\@upc.edu
\\\hline
{} &
~
\\\hline
~
 &
~
\\\hline
~
 &
~
\\\hline
~
 &
~
\\\hline
\end{tabularx}
\end{center}

\bigskip

\begin{center}
\tablefirsthead{}
\tablehead{}
\tabletail{}
\tablelasttail{}
\begin{tabularx}{\textwidth}{|p{2cm}|X|p{2cm}|X|}
\hline
\multicolumn{2}{|l|}{ Written by:} &
\multicolumn{2}{l|}{ Reviewed and approved by:}\\\hline
 Date &
 dd/mm/yyyy &
 Date &
 dd/mm/yyyy\\\hline
 Name &
 Mireia L\'opez Barr\'on &
 Name &
 Alejandro Rodriguez Gomez\\\hline
 Position &
 Project Author &
 Position &
 Project Supervisor\\\hline
\end{tabularx}
\end{center}


%%% ABSTRACT %%%
\clearpage
\newpage
\section*{Abstract}

The need to monitor atmospheric aerosols has grown significantly due to their impact on health and climatology. The Generalized Retrieval of Atmosphere and Surface Properties (GRASP) is a unified algorithm for obtaining these properties from remote sensing observations. This thesis presents the development of a software application that automates the full GRASP execution workflow in quasi-real-time. Built using C\# and the Avalonia UI framework, the cross-platform application seamlessly integrates sequential data processing steps. It automatically downloads photometer data from AERONET and Lidar data from ACTRIS-EARLINET, combines them using Matlab scripts, executes the GRASP algorithm, and visualizes the results. The final product is a reliable, user-friendly tool, fully compatible with Linux environments like the CALCULA system. By automating this complex workflow and incorporating robust error handling, the application significantly reduces manual effort and accelerates atmospheric data analysis.

\input{introduction}
\clearpage

%%% StateOfTheArt %%%
\input{state_of_the_art}
\clearpage

%%% METHODOLOGY %%%
\input{methodology}
\clearpage

%% LIST OF REQUIREMENTS %%%
\input{list_of_requirements}
\clearpage

%%RESEARCH AND PLAN%%% 
\input{research_and_plan}
\clearpage

%%% ARCHITECTURE %%%
\input{architecture}
\clearpage

%%%GUI%%%
\input{gui}
\clearpage

%%%IMPLEMENTATION%%% 
\input{implementation}
\clearpage

%%%TESTING%%% 
\input{testing}
\clearpage

%%% RESULTS %%%
\input{result}
\clearpage

%%%CONCLUSIONS%%% 
\section{Conclusions}
The result of this project is a reliable and easy-to-use application that automatically performs all the steps of the GRASP algorithm almost in real time. After reviewing the whole process from the initial ideas to the final testing, here are the main conclusions:
\smallskip\\\textbf{Main Objectives Achieved:} The main goal was to connect all the steps together in one automatic process, and this has been fully achieved. The application can now easily download data from AERONET and ACTRIS-EARLINET. Then, it runs the Matlab scripts to combine the data, executes the GRASP algorithm, and finally shows the results in a graph. All this automation saves time for the user and makes data analysis much faster.
\smallskip\\\textbf{Working on Different Systems and Handling Errors:} The decision to use C\# and the Avalonia UI package was very important. Thanks to this, the application works perfectly in the CALCULA system, which uses Linux. Also, the program is very strong at handling errors. Because different testing methods were used, such as unit testing and negative testing, the application can deal with problems, like missing data or changed file formats, without crashing. 
\smallskip\\\textbf{Better and Clearer Code:} At the beginning of the project, it was complicated to understand the old code because it was not organized and the file formats were outdated. However, this problem was solved by creating a more logical and structured program. By organizing the code following design patterns and good practices, it was also possible to create unit tests. This means that if someone else wants to work on this project in the future, it will be much easier to understand and improve.
\smallskip\\\textbf{Better User Experience:} The new application is designed specifically for the people who will use it. It has a clear interface that separates the main application settings from the specific project options, making everything simpler. Furthermore, there is an option to use a smaller interface designed for computer screens with a low resolution. This ensures that anyone can use the program, regardless of the computer they have.
\smallskip\\\textbf{Importance of Methodology and Requirements:} Finally, this project has shown how important it is to follow a clear methodology and have a well-defined list of requirements from the start. Because the initial requirements were not very clear, there were many unexpected problems and changes during the development. By eventually defining a strong plan and clear goals, it became much easier to solve problems, organize the code, and finish the application successfully. This proves that spending time on good planning and setting requirements is essential for any software project.

\clearpage

%%% BIBLIOGRAPHY %%%
\newpage

\medskip
\bibliographystyle{unsrt}
\bibliography{bibliography.bib}

%%% ANNEX %%%
\clearpage
\newpage
\appendix
\begin{appendices}
\section{Earlinet Service Class} \label{ap:earlinet_service}
\begin{lstlisting}
public class EarlinetService{
    string baseUrl = "https://api.actris-ares.eu/api/services/restapi/";
    public virtual System.Threading.Tasks.Task<bool> DownloadProductByDateRangeAsync(string type,string FromDate,string ToDate,string outputFolder){
        return DownloadProductByDateRangeAsync(type,FromDate,ToDate,outputFolder,System.Threading.CancellationToken.None);
    }

    public async Task<bool> DownloadProductByDateRangeAsync(string type,string FromDate,string ToDate,string outputFilePath,System.Threading.CancellationToken cancellationToken = default){
        try{
            string url = $"{baseUrl}products/downloads?kind={type}&fromDate={FromDate}&toDate={ToDate}&stations=brc";
            using (HttpClient client = new HttpClient()){
                HttpResponseMessage response = await client.GetAsync(url);
                response.EnsureSuccessStatusCode();
                using (var fs = new FileStream(outputFilePath,FileMode.Create,FileAccess.Write,FileShare.None)){
                    await response.Content.CopyToAsync(fs);
                }
                System.IO.FileInfo f = new FileInfo(outputFilePath);
                if (f.Length <= 1048)
                    return false;
                return true;
            }
        }
        catch (Exception ex){
            return false;
        }
    }
}
\end{lstlisting}
\clearpage
\section{Aeronet Service Class} \label{ap:aeronet_service}
\begin{lstlisting}
public class AeronetService{
    private string baseUrl = "https://aeronet.gsfc.nasa.gov/cgi-bin/";
    public async Task DownloadDataAsync(string destinationFile,string url){
        try{
            using (HttpClient client = new HttpClient()){
                var response = await client.GetAsync(url);
                response.EnsureSuccessStatusCode();

                var content = await response.Content.ReadAsByteArrayAsync();
                await File.WriteAllBytesAsync(destinationFile,content);
            }
        }
        catch (Exception ex){
            Console.WriteLine($"Error downloading data: {ex.Message}");
        }
    }

    public string BuildUrl(DataType _dataType,DateTime startDate,DateTime endDate,string productType = "",string site = "",string product = "",string AVG = "",bool isEnabled = false){
        switch (_dataType){
            case DataType.AerosolInversions:
                if (isEnabled)
                    return BuildUrlAerosolInversions(startDate,endDate,productType,site,product,AVG);
                else
                    return BuildUrlAerosolInversions(startDate,endDate);

            case DataType.OpticalDepth:
                if (isEnabled)
                    return BuildUrlOpticalDepth(startDate,endDate,productType,site,AVG);
                else
                    return BuildUrlOpticalDepth(startDate,endDate,productType);

            case DataType.RawProductsOpticalDepth:
                if (isEnabled)
                    return BuildUrlRawProductsOpticalDepth(startDate,endDate,productType,site,AVG);
                else
                    return BuildUrlRawProductsOpticalDepth(startDate,endDate,productType);
        }
        return "";
    }
    
    {...} //BuildUrl methods can be found in GitHub repository
}
\end{lstlisting}
\clearpage
\section{RunMatlabScript method using ProcessStartInfo class}\label{ap:run_matlab_script}
\begin{lstlisting}
public static void RunMatlabScript (string scriptPath){
    try{
        ProcessStartInfo startInfo = new ProcessStartInfo{
            FileName = "matlab",
            Arguments = $"-batch \"run('{scriptPath}')\"",
            RedirectStandardOutput = true,
            RedirectStandardError = true,
            UseShellExecute = false,
            CreateNoWindow = true
        };
        using var process = new Process { StartInfo = startInfo };
        process.OutputDataReceived += (s,e) =>{
            if (e.Data != null)
                Messenger.Default.Send<string>("WriteMatlabOutput",e.Data);
        };
        process.ErrorDataReceived += (s,e) =>{
            if (e.Data != null)
                Messenger.Default.Send<string>("WriteMatlabErrors",e.Data);
        };
        process.Start();
        process.BeginOutputReadLine();
        process.BeginErrorReadLine();
        process.WaitForExit();
        if (process.ExitCode == 0){ //OK
            Logger.Log($"Script was executed correctly,executing post execution actions . . .");
            script.PostExecutionActions();
        }
        else{
            Logger.Log($"Error occured during execution,executing post execution actions . . .");
            script.PostExecutionActions(false);
        }
    }
    catch (Exception ex){
        Logger.Log($"Error occured during execution,executing post execution actions . . .");
        script.PostExecutionActions(false);
    }
}
\end{lstlisting}
\clearpage
\section{Observable Object}\label{ap:observable_object}
{In order to be able to update the user interface automatically when the data changes, Data Binding is used, which is a feature of the MVVM pattern. It is implemented using ObservableObject as a base class. }
\begin{lstlisting}
public class User : ObservableObject
{
    private string name;

    public string Name
    {
        get => name;
        set => SetProperty(ref name, value);
    }
}
\end{lstlisting}
{
In the example, it can be seen how a variable is updated, when it is detected that the new value is different from the old one, SetProperty is called, which implements INotifyPropertyChanged interface, which is the one in charge to notify the corresponding view. \smallskip\\
In order that the view is updated, it is necessary to call the variable with the same name and to indicate that it is Bindable in the corresponding .axaml file, as it can be seen in the following example:}
\begin{lstlisting}
<TextBlock Text="{Binding Name, Mode=TwoWay, UpdateSourceTrigger=PropertyChanged}" />
\end{lstlisting}
{
In this case, the mode option indicates that the variable can be updated from the view to the model and vice versa, and the update source trigger option indicates that the update will be done when the property changes. Other options can be used, but these are the ones that are most common and that have been used in this project.
}

\clearpage
\section{Relay Command}\label{ap:relay_command}
{Relay command is used to implement commands in the MVVM pattern. It is implemented using ICommand as a base class. It can be used in different ways, but in this project it has been implemented using the following pattern for all View Models:}
\begin{lstlisting}
public class ViewModel : ObservableObject
{
  public ICommand Cmd => new RelayCommand(Execute, CanExecute);
  private void Execute(object _)
  {
      ExecutionActions();
  }

  private bool CanExecute(object _)
  {
      return true;
  }
}
\end{lstlisting}
{In order to execute the command, it is necessary to link an interface element of the .axaml file to the Execute method}
\begin{lstlisting}
<Button Content="Execute" Command="{Binding Cmd}" />
\end{lstlisting}

\clearpage
\section{Messenger}\label{ap:messenger}
The use of Messenger class consists on defining a Send Action, that can include some parameter or variable and a Receive Action, that will be executed when the Send Action is called, executing the linked method.\smallskip\\ 
In the following example it can be seen an example of how the DataCombinationViewModel is using the Messenger class to receive modifications in the buttons enabled state:\smallskip
\begin{lstlisting}
public DataCombinationViewModel()
{
    Messenger.Default.Register<bool>("UpdateButtonsEnabled", UpdateButtonsEnabled);
}   

private void UpdateButtonsEnabled(bool status)
{
    AreButtonsEnabled = status;
    if (!status)
    {
        isEnabled_D1_L = false;
        isEnabled_D1P_L = false;
        isEnabled_D1_L_VD = false;
        isEnabled_D1P_L_VD = false;
    }
}
\end{lstlisting}
In the corresponding controller, in order to send the modification in order to enable the buttons, for example, it is necessary to execute the following line:
\begin{lstlisting}
Messenger.Default.Send<bool>("UpdateButtonsEnabled", true);
\end{lstlisting}
\section{Test Results} \label{ap:test_results}
\includepdf[pages=-]{"./Tests.pdf"}

\section{User Manual} \label{ap:user_manual}
\includepdf[pages=-]{"../User Manual/main.pdf"}

\end{appendices}

\end{document}
